% Introduction and the first two Results sections. Numbers come from numbers.tex (generated).

\noindent Sentiment analysis turns a piece of writing into a number that says how positive or negative it is. The number is computed for product reviews, posts, news, clinical notes and parliamentary speech, often for millions of texts at a time\cite{pang2002,dodds2011,golder2011,proksch2019,ribeiro2016}. Two families of method are in use. Lexicon methods look each word up in a list of human word ratings and combine the values by rule\cite{dodds2010,hutto2014,taboada2011,tausczik2010}. Classifiers, now mostly fine-tuned transformers, are trained on labelled text and return a class probability\cite{devlin2019,liu2019roberta,sanh2019}.

Both families return a bounded number for each text: a probability, a mean, or a sum squashed into the interval from $-1$ to $1$. A bounded score is intensive in the sense that temperature is intensive. It does not grow with the amount of text, and the value for a whole cannot be computed from the values of its parts unless the counts behind them are kept; a class probability and a squashed sum have none to keep. Three practical problems follow. A transformer of the BERT family reads at most 512 tokens at a time\cite{devlin2019}, so a long document is cut into chunks whose scores are averaged, and the average depends on where the cuts fall. A squashed sum saturates: in \vaderBooks{} of the novels studied below, VADER's compound score is beyond $\pm0.99$ for \vaderSaturated\% of whole chapters. And an average of averages can rank two bodies of text in the opposite order to their pooled values, the reversal known from Simpson's paradox\cite{simpson1951,blyth1972,lerman2018}. The third problem is a risk and not an observation: in the data below the two orderings disagree for few pairs, and no comparison reverses.

Here I treat sentiment as a quantity that adds. Every word carries a positive and a negative charge taken from human word ratings. A span of text is summarised by a short vector, its signature, which holds counts and sums of charge weighted by powers of word position. Signatures of adjacent spans merge exactly into the signature of the joined span, so a novel is scored in one pass and every chapter, paragraph and sentence can be read off afterwards without scoring anything again. Scores that people can compare, such as net charge per thousand words or a predicted human rating with an error bar, are read from the signature at whatever scale is wanted.

The parts have precedents. The hedonometer averages human word ratings\cite{dodds2010,dodds2011}. Word-count tools and studies of mood keep positive and negative words apart\cite{tausczik2010,thelwall2010,golder2011}, in line with evidence that the two are separable in evaluation\cite{cacioppo1994}. Political scientists score texts from counts of opposed items on a log-ratio scale with a stated uncertainty\cite{lowe2011,proksch2019}. Summaries that merge exactly across partitions, including sums weighted by polynomial or exponential functions of time, are standard in stream processing\cite{cormode2009,pebay2008,agarwal2013,boykin2014}. A model that composes the sentiment of a phrase by multiplying matrices assigned to its words is associative in the same way\cite{yessenalina2011}. Emotional arcs of stories have been drawn with moving windows, transforms and fitted polynomials\cite{reagan2016,jockers2015,bizzoni2023}. This paper joins these into one object for sentiment, states which position weights such an object can carry, and tests it against human ratings.

The aim is a mathematical measure of sentiment for which granularity holds: the score of a whole is the merge of the scores of its parts, at every level from the sentence to the book. The work began with an unpublished draft I wrote in 2024 (its scripts are in the repository), which proposed that feeling in text could be scored as a number from word values, a decay over position and a penalty for words that carry no feeling. It made two claims about how people write and read. Front-loading: words that carry feeling come early in a text and count for more there. Dilution: words that carry none weaken the feeling of the whole. I wrote both down as hypotheses with decision rules, and committed them to the project repository, before scoring any evaluation data (Methods).

Front-loading fails as a claim about placement. In sentences, tweets and posts charged words gather at the end, and in reviews and novel chapters they are spread almost evenly; nowhere do they gather at the start. As a claim about weight it depends on the genre. Readers of film-review sentences weigh the early part more, and readers of tweets and of Yelp reviews the late part. Midway through the work I replaced front-loading with a hypothesis of recency, that what comes last counts most, and committed three tests of it. All three passed by their stated rules. Two of them compared parts of unequal length, and when a referee agent pointed this out and I held length fixed, both reversed. I report the sequence in full and withdraw the hypothesis. Dilution holds, by the planned test in \winsNoDil{} of six corpora and in human phrase ratings on either side of the feeling. A last experiment applies the measure to a current question and finds that essays written by language models differ from student essays on the same assignments in how they carry feeling.

The code, the analyses and the text of this paper were produced with an AI assistant working under my direction, and AI referee agents reviewed the draft (Methods). ``I'' is the author, who set the aims and the hypotheses and answers for the claims.

\section{Results}

\subsection{A signature that merges}

Take a text as a sequence of words. Word $i$ has a positive charge $h_i$ and a negative charge $d_i$, each between 0 and 1, at most one of them nonzero. The values come from two open sets of human word ratings, the labMT list\cite{dodds2011} (\nLabmt{} words, 50 raters each) and the VADER list\cite{hutto2014} (\nVader{} entries, 10 raters each). On the \nShared{} words they share, the two sets correlate at $r=\rLexAgree$. I put both on one scale, on which 1 is the positive end of the rating scale and 0 its neutral point, and average them (Methods). Positive charge is measured in hedons (hd) and negative charge in dolors (dl), the names that moral philosophy uses for units of pleasure and pain\cite{mackaye1906,feldman1978}. One hedon is the charge of a word that every rater places at the top of the scale; \emph{love} carries \hdLove\hd{} and \emph{terrible} \dlTerrible\dl. A word rated within a quarter of the scale from neutral gets no charge and is called a buffer word, following the neutral band of the hedonometer\cite{dodds2011}. In the six rated corpora used here, \chargedShare\% of running words are charged.

The signature of a span holds the number of words $L$, the number of charged words $E$, and the sums
\begin{equation}
S^{+}_{k}=\sum_{i} h_i\, i^{k},\qquad S^{-}_{k}=\sum_{i} d_i\, i^{k},\qquad k=0,\dots,3,
\end{equation}
where $i$ is the position of the word in the span, counted from zero. $S^{+}_{0}$ and $S^{-}_{0}$ are the span's total hedons and dolors. The signature also holds the hedons and dolors that fall under a negation, after an intensifier, after a diminisher, and after a contrast word such as \emph{but}, together with the sum of squared charges (Methods). That is 19 numbers, whatever the length of the span (\figref{fig:schematic}a,b).

Signatures merge. If span $A$ ends at a sentence end and is followed by span $B$, the counts and totals add, and the position sums of $B$ are shifted by the length of $A$:
\begin{equation}
S_{k}(AB)=S_{k}(A)+\sum_{j=0}^{k}\binom{k}{j}L_A^{\,k-j}\,S_{j}(B).
\label{eq:merge}
\end{equation}
The operation is associative and can be undone, so the signature of $B$ follows from those of $AB$ and $A$ (Proposition~1, Methods). At a cut inside a sentence the counts, the totals and the position sums still merge exactly, and the operator entries can differ, because the reach of an operator may cross the cut. A text of any length is therefore scored in one pass with constant memory. Parts can be scored separately and merged in reading order to the same result as a single pass. A table of prefix signatures, one row per sentence, gives the signature of any contiguous run of sentences from two of its rows. In integer arithmetic these statements hold to the last digit, and the reference implementation, which counts in thousandths of a hedon, is tested for exact equality on random texts. The fast implementation uses double-precision floats; over \nNovels{} novels its merged and single-pass signatures agree to a relative error below \mergeErr. In floats the two-row query is accurate at any depth for counts and totals, and so for density and level. The position sums of a short run deep in a long text are lost to cancellation, so the arc or the predicted rating of such a run is read by merging the run's own rows, at a cost that grows with its length.

Not every position weight allows a merge. Suppose a summary keeps finitely many sums $\sum_i c_i\,\varphi(i)$ of charge against weight functions $\varphi$. The summary of $AB$ can be computed from the summaries of $A$ and $B$ and their lengths exactly when the weight functions span a space that is closed under shifting the position, and such weights are combinations of $i^{k}\rho^{i}$: polynomials, exponentials and their products (Proposition~2, Methods). The statement is an instance of a classical fact about translation-invariant function spaces\cite{anselone1964,szekelyhidi1982} and about functions on strings with finite-dimensional representations\cite{schutzenberger1961,berstel2010}, and stream-processing systems rely on it\cite{cormode2009}. The decay factor of the 2024 draft, a weight that falls geometrically with position, is therefore one of the few position weights that survive a merge, at the cost of one number per decay rate. And among mergeable weights, polynomials are the only ones that give a continuous function of relative position, that is, of position divided by a length that is unknown until the text ends (Methods). The signature keeps polynomial sums for that reason: the sum of charge under any cubic function of relative position can be read from it, for a post and for a novel alike.

\subsection{Reading a score from a signature}

A signature is extensive. Scores that can be compared between texts are read from it (\figref{fig:schematic}c).

\emph{Totals} are the hedons $H=S^{+}_{0}$ and dolors $D=S^{-}_{0}$ over $L$ words. \emph{Density} is the net charge per thousand words, $1000\,(H-D)/L$; the hedonometer's average happiness is a density taken over charged words only.

\emph{Level} follows a model from the psychology of judgement. In Anderson's averaging model a judgement is a weighted mean of the pieces of information received and an initial neutral impression\cite{anderson1965}, and information that says nothing pulls a judgement toward neutral, the dilution effect\cite{nisbett1981}. For words this gives
\begin{equation}
\text{level}=1000\,\frac{H-D}{a+E+\kappa B},
\label{eq:level}
\end{equation}
where $B=L-E$ counts buffer words, $a$ is the weight of the neutral first impression in units of charged words, and $\kappa$ is the weight of a buffer word relative to a charged one. With $\kappa=0$ buffer words do nothing; with $a=0$ and $\kappa=1$ the level is the density. The denominator is the exposure, the amount of evidence behind the score, and it gives the level a standard error that shrinks as a text grows (Methods).

\emph{Calibrated valence} is a predicted human rating. It is a linear combination of ten signature entries, each divided by the exposure: hedons, dolors, three position moments of net charge, negated hedons, negated dolors, and the net charge that is intensified, diminished, and after a contrast word. Given $a$ and $\kappa$, the coefficients solve a ridge regression against human ratings, a convex problem with a closed-form solution; $a$ and $\kappa$ are chosen on a grid (Methods). The fitted coefficients say how much a negation flips a word and how position is weighed; they are estimates with intervals, in the tradition of rule constants measured from ratings\cite{hutto2014,kiritchenko2016}.

The \emph{arc} is a cubic fit of net charge against relative position. It has a closed form in the position sums and equals the least-squares cubic up to terms that shrink as the square of the span's length (Methods), so it costs nothing beyond the signature and suits spans of a few dozen words or more.

